% Front matter for the near-root Version 4 verification manuscript.
% The compile-time switch \ifErdosProofClosed is defined by the master file.

\begin{abstract}
\ifErdosProofClosed
Let $G_n\sim G(n,1/2)$. We prove, along the full sequence of integers, that
\[
  \chi(G_n)-\zeta(G_n)
  \ge
  \frac{(\log 2)^2}{4}
  \log\!\left(\frac{1000}{639}\right)
  \frac{n}{(\log n)^3}
\]
with probability tending to one. The proof compares the ordinary coloring
root with a signed four-size cocoloring root. A one-part buffer above the signed
root gives a first-moment margin of order $(\log n)^2$, which is already enough
for the full-corner second-moment estimate and therefore retains the complete
leading root separation. Exact sign compatibility, completion-free high-cell
fibers, a fused endpoint kernel, and a critical-quarter residual split give a
seed exponent $O(\sqrt n(\log n)^{3/2})$; bounded differences then amplify the
seed at negligible cost.
\else
Let $G_n\sim G(n,1/2)$. This verification draft gives a self-contained
candidate derivation of a full-sequence lower bound of order
$n/(\log n)^3$ for $\chi(G_n)-\zeta(G_n)$, conditional on the exact review and
formal-replay obligations stated in the appendices. Its new point is that
midpoint placement is unnecessary: the integer
$\lceil r_4^{\mathrm{co}}\rceil+1$ has signed first moment
$\exp(\Theta((\log n)^2))$, which dominates the only polynomial full-corner
multiplicity and retains the full first-moment root separation. The manuscript
keeps the publication switch disabled until this one-part-buffer interface and
the existing phase, partial-diagonal, skeleton, and attachment packages are
replayed together.
\fi
\end{abstract}

\maketitle

\ifErdosProofClosed\else
\begin{center}
\small\emph{Verification status.} The one-part-buffer placement and doubled
coefficient are complete candidate paper interfaces, not yet independently
certified. The theorem remains conditional until the exact declarations,
private Lean replay, and integrated dependency audit listed in the appendices
have all passed on one commit.
\end{center}
\fi

\section*{Introduction}
\addcontentsline{toc}{section}{Introduction}
\label{sec:introduction-near-root-v4}

A \emph{cocoloring} of a graph $G$ is a partition of $V(G)$ into nonempty
classes, each inducing either an empty graph or a complete graph. The least
number of classes in such a partition is the \emph{cochromatic number}
$\zeta(G)$. Since an ordinary proper coloring is a cocoloring using only empty
classes, $\zeta(G)\le\chi(G)$.

Let $G_n\sim G(n,1/2)$ be the labeled random graph on $[n]$. Erd\H{o}s and
Gimbel asked whether
\[
  \chi(G_n)-\zeta(G_n)\longrightarrow\infty
\]
with high probability \citep[p.~263]{erdos-gimbel-1993}. The question was
restated by Gimbel \citep[Section~7.4]{gimbel-2016} and is cataloged as
Erd\H{o}s Problem~625 \citep{bloom-erdos625}.

\ifErdosProofClosed
\begin{maintheorem}
For $G_n\sim G(n,1/2)$,
\[
  \mathbb P\!\left(
    \chi(G_n)-\zeta(G_n)
    \ge
    \frac{(\log 2)^2}{4}
    \log\!\left(\frac{1000}{639}\right)
    \frac{n}{(\log n)^3}
  \right)
  \longrightarrow1.
\]
In particular, the difference tends to infinity with high probability along
the full sequence of integers.
\end{maintheorem}
\else
\begin{maintheorem}[Conditional verification statement]
Assume the phase and finite-dual package, the one-part-buffer signed profile,
the complete partial-diagonal estimate, and the global
skeleton--attachment assembly listed in the appendices. Then
\[
  \mathbb P\!\left(
    \chi(G_n)-\zeta(G_n)
    \ge
    \frac{(\log 2)^2}{4}
    \log\!\left(\frac{1000}{639}\right)
    \frac{n}{(\log n)^3}
  \right)
  \longrightarrow1.
\]
\end{maintheorem}
\fi

The full-sequence quantifier is substantive because the natural class-size
cutoff jumps when the independence-number phase crosses an integer. Every
asymptotic estimate below is uniform up to both endpoints of each phase
interval. The phase, finite-dual, target-transport, partial-diagonal, endpoint,
and residual errors are named separately so that the same overlap feature is
not charged twice.

\subsection*{Relation to previous work}

The first-order asymptotic for the chromatic number of a dense random graph was
proved by Bollob\'as \citep{bollobas-1988}, following the early work of
Grimmett and McDiarmid \citep{grimmett-mcdiarmid-1975}; later refinements
include \citet{mcdiarmid-1990}, \citet{panagiotou-steger-2009}, and
\citet{heckel-2018}. The cochromatic number belongs to the theory of generalized
chromatic numbers associated with hereditary graph properties
\citep{scheinerman-1992,bollobas-thomason-1995}.

For the difference $\chi(G_n)-\zeta(G_n)$, phase-dependent progress was
obtained through the nonconcentration of the chromatic number
\citep{heckel-2024-question,steiner-2024,heckel-2025-difference}. The present
candidate proof uses a four-size signed profile that remains valid through the
complete phase, evaluates sign compatibility exactly, and carries the rare
second-moment seed to high probability by a quantitative exposure argument.

\subsection*{Why one extra class is enough}

Let $r_+(n)$ be the ordinary first-moment root and
$r_4^{\mathrm{co}}(n)$ the signed four-size root. Their separation is
\[
  r_+-r_4^{\mathrm{co}}
  =
  \left[
    \frac{(\log 2)^2}{4}A_4(\delta_n)+o(1)
  \right]
  \frac{n}{(\log n)^3},
\]
where $A_4(\delta)=\log 2-D_4(\delta)$ is uniformly positive. The earlier
verification draft placed the signed profile at the midpoint and consequently
retained one half of this separation.

The stronger placement is
\[
  k_{\mathrm{co}}
  =\left\lceil r_4^{\mathrm{co}}\right\rceil+1.
\]
Its distance from the signed root lies in $[1,2)$. Since the derivative of the
signed objective is uniformly of order $(\log n)^2$, the resulting exact
integer profile has
\[
  \log\mathbb E Z_{\mathbf k}^{\mathrm{sgn}}
  =\Theta((\log n)^2).
\]
This is smaller than the midpoint margin, but it is sufficient. The central
partial-diagonal range needs only
$k_{\mathrm{co}}^{-1}\log\mathbb E Z=o(1)$, and the full corner has only
$(k_{\mathrm{co}}+1)^4$ residual profiles. Hence
\[
  (k_{\mathrm{co}}+1)^4
  \bigl(\mathbb E Z_{\mathbf k}^{\mathrm{sgn}}\bigr)^{-1}
  =o(1).
\]
No other part of the second-moment argument consumes a macroscopic distance
from the signed root. The complete leading root separation is therefore
available in the final gap.

\subsection*{Proof components}

The proof has four components.

\begin{enumerate}
\item \textbf{Uniform root geometry.} The finite support comparison is exact
before passage to the limiting Gaussian dual. The signed profile is then
placed one integer part above the signed root and tangent-rounded while
preserving both conservation laws.

\item \textbf{Exact signed overlap.} For two signed witnesses, compatibility
factors into local cell rewards and the cardinality of a binary cycle space.

\item \textbf{Fused high-cell transport.} High cells form a matching. Their
complete physical fibers are summed before estimates are taken. Deficit
summation and endpoint transport are fused into one nonnegative $4\times4$
kernel, giving a logarithmic skeleton cost
$O(\sqrt n(\log n)^{3/2})$.

\item \textbf{Critical-quarter attachments and amplification.} A fixed-buffer
threshold at the intrinsic scale $2^{U/4}$ bounds the residual attachment by
$\exp\{O(\sqrt n(\log n)^{3/2})\}$. Paley--Zygmund supplies a rare witness,
and bounded differences amplify it while adding only
$O(n^{3/4}(\log n)^{-1/4})=o(n/(\log n)^3)$ classes.
\end{enumerate}

The exact four-support ledger proves
\[
  A_4(\delta)
  >
  \log\!\left(\frac{20000}{12777}\right)
  =
  \log\!\left(\frac{1000}{639}\right)
  +
  \log\!\left(\frac{12780}{12777}\right).
\]
The last term is fixed slack. Retaining the full root separation gives the
coefficient
\[
  \frac{(\log 2)^2}{4}
  \log\!\left(\frac{1000}{639}\right)
  =0.053792819616758\ldots.
\]

Sections~1--5 establish the phase and one-part-buffer profile. Section~6 gives
the exact signed-overlap identity, and Section~7 proves the partial-diagonal
estimate with the revised full corner. Sections~8--9 contain the fused endpoint
and critical-quarter residual bounds. Section~10 proves the amplifier, the
final section performs the coefficient ledger, and the appendices record the
paper-to-Lean interfaces and remaining gates.
